% GENERATED FILE - do not hand-edit.
% Source of truth: data/calculus-reviewed/c-1.json in the zenith-website repo.
% Regenerate with: python3 scripts/build_calculus_tex.py
\documentclass[11pt]{article}
\usepackage[utf8]{inputenc}
\usepackage[T1]{fontenc}
\usepackage{amsmath,amssymb}
\usepackage{graphicx}
\usepackage[margin=1in]{geometry}
\graphicspath{{../..//}}
\setlength{\parindent}{0pt}
\setlength{\parskip}{0.6em}
\begin{document}
\section*{Chapter 1: Functions --- Practice Exercises C1}

\bigskip\noindent\textbf{A1.}\ If $f(x)=x^3-2x-1$, then $f(-2)=$
\begin{list}{}{\setlength{\leftmargin}{3em}\setlength{\labelwidth}{2.5em}\setlength{\itemsep}{0pt}\setlength{\topsep}{2pt}}
\item[(A)] $-13$
\item[(B)] $-5$
\item[(C)] $-1$
\item[(D)] $7$
\end{list}

\bigskip\noindent\textbf{A2.}\ The domain of $f(x)=\dfrac{x-1}{x^2+1}$ is
\begin{list}{}{\setlength{\leftmargin}{3em}\setlength{\labelwidth}{2.5em}\setlength{\itemsep}{0pt}\setlength{\topsep}{2pt}}
\item[(A)] all $x\ne1$
\item[(B)] all $x\ne1,-1$
\item[(C)] all $x\ne-1$
\item[(D)] all reals
\end{list}

\bigskip\noindent\textbf{A3.}\ The domain of $g(x)=\dfrac{\sqrt{x-2}}{x^2-x}$ is
\begin{list}{}{\setlength{\leftmargin}{3em}\setlength{\labelwidth}{2.5em}\setlength{\itemsep}{0pt}\setlength{\topsep}{2pt}}
\item[(A)] all $x\ne0,1$
\item[(B)] $x\le2,\ x\ne0,1$
\item[(C)] $x\ge2$
\item[(D)] $x>2$
\end{list}

\bigskip\noindent\textbf{A4.}\ If $f(x)=x^3-3x^2-2x+5$ and $g(x)=2$, then $g(f(x))=$
\begin{list}{}{\setlength{\leftmargin}{3em}\setlength{\labelwidth}{2.5em}\setlength{\itemsep}{0pt}\setlength{\topsep}{2pt}}
\item[(A)] $2x^3-6x^2-2x+10$
\item[(B)] $2x^2-6x+1$
\item[(C)] $-3$
\item[(D)] $2$
\end{list}

\bigskip\noindent\textbf{A5.}\ If $f(x)=x^3-3x^2-2x+5$ and $g(x)=2$, then $f(g(x))=$
\begin{list}{}{\setlength{\leftmargin}{3em}\setlength{\labelwidth}{2.5em}\setlength{\itemsep}{0pt}\setlength{\topsep}{2pt}}
\item[(A)] $2x^3-6x^2-2x+10$
\item[(B)] $2x^2-6x+1$
\item[(C)] $-3$
\item[(D)] $2$
\end{list}

\bigskip\noindent\textbf{A6.}\ If $f(x)=x^3+Ax^2+Bx-3$, $f(1)=4$, and $f(-1)=-6$, what is $2A+B$?
\begin{list}{}{\setlength{\leftmargin}{3em}\setlength{\labelwidth}{2.5em}\setlength{\itemsep}{0pt}\setlength{\topsep}{2pt}}
\item[(A)] $12$
\item[(B)] $8$
\item[(C)] $0$
\item[(D)] $-2$
\end{list}

\bigskip\noindent\textbf{A7.}\ Which equation has a graph symmetric with respect to the origin?
\begin{list}{}{\setlength{\leftmargin}{3em}\setlength{\labelwidth}{2.5em}\setlength{\itemsep}{0pt}\setlength{\topsep}{2pt}}
\item[(A)] $y=\dfrac{x-1}{x}$
\item[(B)] $y=2x^4+1$
\item[(C)] $y=x^3+2x$
\item[(D)] $y=x^3+2$
\end{list}

\bigskip\noindent\textbf{A8.}\ Let $g$ be defined for all reals. Which condition is not sufficient to guarantee that $g$ has an inverse function?
\begin{list}{}{\setlength{\leftmargin}{3em}\setlength{\labelwidth}{2.5em}\setlength{\itemsep}{0pt}\setlength{\topsep}{2pt}}
\item[(A)] $g(x)=ax+b,\ a\ne0$
\item[(B)] $g$ is strictly decreasing
\item[(C)] $g$ is symmetric to the origin
\item[(D)] $g$ is one-to-one
\end{list}

\bigskip\noindent\textbf{A9.}\ Let $f(x)=\sin(\arctan x)$. The range of $f$ is
\begin{list}{}{\setlength{\leftmargin}{3em}\setlength{\labelwidth}{2.5em}\setlength{\itemsep}{0pt}\setlength{\topsep}{2pt}}
\item[(A)] $\{y\mid-1<y<1\}$
\item[(B)] $\{y\mid-1\le y\le1\}$
\item[(C)] $\{y\mid-\pi/2<y<\pi/2\}$
\item[(D)] $\{y\mid-\pi/2\le y\le\pi/2\}$
\end{list}

\bigskip\noindent\textbf{A10.}\ Let $g(x)=|\cos x-1|$. At which $x$ does $g$ attain its maximum on $[0,2\pi]$?
\begin{list}{}{\setlength{\leftmargin}{3em}\setlength{\labelwidth}{2.5em}\setlength{\itemsep}{0pt}\setlength{\topsep}{2pt}}
\item[(A)] $0$
\item[(B)] $\pi/2$
\item[(C)] $2$
\item[(D)] $\pi$
\end{list}

\bigskip\noindent\textbf{A11.}\ Which function is not odd?
\begin{list}{}{\setlength{\leftmargin}{3em}\setlength{\labelwidth}{2.5em}\setlength{\itemsep}{0pt}\setlength{\topsep}{2pt}}
\item[(A)] $f(x)=\sin x$
\item[(B)] $f(x)=\sin2x$
\item[(C)] $f(x)=x^3+1$
\item[(D)] $f(x)=\dfrac{x}{x^2+1}$
\end{list}

\bigskip\noindent\textbf{A12.}\ If the solutions of $f(x)=0$ are $x=1,-2$, then $f(2x)=0$ at
\begin{list}{}{\setlength{\leftmargin}{3em}\setlength{\labelwidth}{2.5em}\setlength{\itemsep}{0pt}\setlength{\topsep}{2pt}}
\item[(A)] $x=\frac12,-1$
\item[(B)] $x=-\frac12,1$
\item[(C)] $x=2,-4$
\item[(D)] $x=-2,4$
\end{list}

\bigskip\noindent\textbf{A13.}\ The set of zeros of $f(x)=x^3+4x^2+4x$ is
\begin{list}{}{\setlength{\leftmargin}{3em}\setlength{\labelwidth}{2.5em}\setlength{\itemsep}{0pt}\setlength{\topsep}{2pt}}
\item[(A)] $\{-2\}$
\item[(B)] $\{0,-2\}$
\item[(C)] $\{0,2\}$
\item[(D)] $\{2\}$
\end{list}

\bigskip\noindent\textbf{A14.}\ At which $x$ do $y=x+2$ and $y^2=4x$ intersect?
\begin{list}{}{\setlength{\leftmargin}{3em}\setlength{\labelwidth}{2.5em}\setlength{\itemsep}{0pt}\setlength{\topsep}{2pt}}
\item[(A)] $-2$ and $2$
\item[(B)] $-2$
\item[(C)] $2$
\item[(D)] no intersection
\end{list}

\bigskip\noindent\textbf{A15.}\ Reflect $f(x)=1-3^x$ in the $y$-axis. The resulting function is
\begin{list}{}{\setlength{\leftmargin}{3em}\setlength{\labelwidth}{2.5em}\setlength{\itemsep}{0pt}\setlength{\topsep}{2pt}}
\item[(A)] $g(x)=1-3^{-x}$
\item[(B)] $g(x)=3^x-1$
\item[(C)] $g(x)=\log_3(x-1)$
\item[(D)] $g(x)=\log_3(1-x)$
\end{list}

\bigskip\noindent\textbf{A16.}\ Let $g$ be the inverse function of $f$. Then $f(g(x))=$
\begin{list}{}{\setlength{\leftmargin}{3em}\setlength{\labelwidth}{2.5em}\setlength{\itemsep}{0pt}\setlength{\topsep}{2pt}}
\item[(A)] $1$
\item[(B)] $x$
\item[(C)] $1/x$
\item[(D)] $f(x)g(x)$
\end{list}

\bigskip\noindent\textbf{A17.}\ The function $f(x)=2x^3+x-5$ has exactly one real zero. It is between
\begin{list}{}{\setlength{\leftmargin}{3em}\setlength{\labelwidth}{2.5em}\setlength{\itemsep}{0pt}\setlength{\topsep}{2pt}}
\item[(A)] $-1$ and $0$
\item[(B)] $0$ and $1$
\item[(C)] $1$ and $2$
\item[(D)] $2$ and $3$
\end{list}

\bigskip\noindent\textbf{A18.}\ The period of $f(x)=\sin\left(\dfrac{2\pi}{3}x\right)$ is
\begin{list}{}{\setlength{\leftmargin}{3em}\setlength{\labelwidth}{2.5em}\setlength{\itemsep}{0pt}\setlength{\topsep}{2pt}}
\item[(A)] $2/3$
\item[(B)] $3/2$
\item[(C)] $3$
\item[(D)] $6$
\end{list}

\bigskip\noindent\textbf{A19.}\ The range of $f(x)=\ln(\cos x)$ is
\begin{list}{}{\setlength{\leftmargin}{3em}\setlength{\labelwidth}{2.5em}\setlength{\itemsep}{0pt}\setlength{\topsep}{2pt}}
\item[(A)] $\{y\mid-\infty<y\le0\}$
\item[(B)] $\{y\mid0<y\le1\}$
\item[(C)] $\{y\mid-1<y<1\}$
\item[(D)] $\{y\mid0\le y\le1\}$
\end{list}

\bigskip\noindent\textbf{A20.}\ If $\log_b(3^b)=b/2$, then $b=$
\begin{list}{}{\setlength{\leftmargin}{3em}\setlength{\labelwidth}{2.5em}\setlength{\itemsep}{0pt}\setlength{\topsep}{2pt}}
\item[(A)] $1/9$
\item[(B)] $1/2$
\item[(C)] $3$
\item[(D)] $9$
\end{list}

\bigskip\noindent\textbf{A21.}\ The inverse of $f(x)=x^3+2$ is
\begin{list}{}{\setlength{\leftmargin}{3em}\setlength{\labelwidth}{2.5em}\setlength{\itemsep}{0pt}\setlength{\topsep}{2pt}}
\item[(A)] $\dfrac1{x^3-2}$
\item[(B)] $(x-2)^3$
\item[(C)] $\sqrt[3]{x+2}$
\item[(D)] $\sqrt[3]{x-2}$
\end{list}

\bigskip\noindent\textbf{A22.}\ The set of $x$-intercepts of $f(x)=x^3-2x^2-x+2$ is
\begin{list}{}{\setlength{\leftmargin}{3em}\setlength{\labelwidth}{2.5em}\setlength{\itemsep}{0pt}\setlength{\topsep}{2pt}}
\item[(A)] $\{-1,1\}$
\item[(B)] $\{1,2\}$
\item[(C)] $\{-1,1,2\}$
\item[(D)] $\{-1,-2,2\}$
\end{list}

\bigskip\noindent\textbf{A23.}\ On $(-\pi/2,\pi/2)$, the range of $f(x)=e^{\tan x}$ is
\begin{list}{}{\setlength{\leftmargin}{3em}\setlength{\labelwidth}{2.5em}\setlength{\itemsep}{0pt}\setlength{\topsep}{2pt}}
\item[(A)] all reals
\item[(B)] positive reals
\item[(C)] nonnegative reals
\item[(D)] $\{y\mid0<y\le1\}$
\end{list}

\bigskip\noindent\textbf{A24.}\ Which is a reflection of $y=f(x)$ in the $x$-axis?
\begin{list}{}{\setlength{\leftmargin}{3em}\setlength{\labelwidth}{2.5em}\setlength{\itemsep}{0pt}\setlength{\topsep}{2pt}}
\item[(A)] $y=-f(x)$
\item[(B)] $y=f(-x)$
\item[(C)] $y=f(|x|)$
\item[(D)] $y=-f(-x)$
\end{list}

\bigskip\noindent\textbf{A25.}\ The smallest positive $x$ for which $f(x)=\sin(x/3)-1$ is a maximum is
\begin{list}{}{\setlength{\leftmargin}{3em}\setlength{\labelwidth}{2.5em}\setlength{\itemsep}{0pt}\setlength{\topsep}{2pt}}
\item[(A)] $\pi/2$
\item[(B)] $\pi$
\item[(C)] $3\pi/2$
\item[(D)] $3\pi$
\end{list}

\bigskip\noindent\textbf{A26.}\ $\tan\left(\arccos\left(-\dfrac{\sqrt2}{2}\right)\right)=$
\begin{list}{}{\setlength{\leftmargin}{3em}\setlength{\labelwidth}{2.5em}\setlength{\itemsep}{0pt}\setlength{\topsep}{2pt}}
\item[(A)] $-1$
\item[(B)] $-\sqrt3/3$
\item[(C)] $\sqrt3/3$
\item[(D)] $1$
\end{list}

\bigskip\noindent\textbf{A27.}\ The inverse of $f(x)=2e^{-x}$ is
\begin{list}{}{\setlength{\leftmargin}{3em}\setlength{\labelwidth}{2.5em}\setlength{\itemsep}{0pt}\setlength{\topsep}{2pt}}
\item[(A)] $\ln(2/x)$
\item[(B)] $\ln(x/2)$
\item[(C)] $\frac12\ln x$
\item[(D)] $\sqrt{\ln x}$
\end{list}

\bigskip\noindent\textbf{A28.}\ Which function does not have an inverse function?
\begin{list}{}{\setlength{\leftmargin}{3em}\setlength{\labelwidth}{2.5em}\setlength{\itemsep}{0pt}\setlength{\topsep}{2pt}}
\item[(A)] $y=\sin x,\ -\pi/2\le x\le\pi/2$
\item[(B)] $y=x^3+2$
\item[(C)] $y=\dfrac{x}{x^2+1}$
\item[(D)] $y=\ln(x-2),\ x>2$
\end{list}

\bigskip\noindent\textbf{A29.}\ Let $f(x)=\ln x$ and $g(x)=9-x^2$. The domain of $f(g(x))$ is
\begin{list}{}{\setlength{\leftmargin}{3em}\setlength{\labelwidth}{2.5em}\setlength{\itemsep}{0pt}\setlength{\topsep}{2pt}}
\item[(A)] $\{x\mid x\le3\}$
\item[(B)] $\{x\mid |x|>3\}$
\item[(C)] $\{x\mid |x|<3\}$
\item[(D)] $\{x\mid0<x<3\}$
\end{list}

\bigskip\noindent\textbf{A30.}\ Let $f(x)=\ln x$ and $g(x)=9-x^2$. The range of $f(g(x))$ is
\begin{list}{}{\setlength{\leftmargin}{3em}\setlength{\labelwidth}{2.5em}\setlength{\itemsep}{0pt}\setlength{\topsep}{2pt}}
\item[(A)] $\{y\mid y>0\}$
\item[(B)] $\{y\mid0<y\le\ln9\}$
\item[(C)] $\{y\mid y\le\ln9\}$
\item[(D)] $\{y\mid y<0\}$
\end{list}

\bigskip\noindent\textbf{A31.}\ The curve $x(t)=t^2+3$, $y(t)=t^2+4$ is part of a(n)
\begin{list}{}{\setlength{\leftmargin}{3em}\setlength{\labelwidth}{2.5em}\setlength{\itemsep}{0pt}\setlength{\topsep}{2pt}}
\item[(A)] line
\item[(B)] circle
\item[(C)] parabola
\item[(D)] ellipse
\end{list}

\bigskip\noindent\textbf{A32.}\ Which equation includes $x(t)=\cos^2t$, $y(t)=2\sin t$?
\begin{list}{}{\setlength{\leftmargin}{3em}\setlength{\labelwidth}{2.5em}\setlength{\itemsep}{0pt}\setlength{\topsep}{2pt}}
\item[(A)] $x^2+y^2=4$
\item[(B)] $4x^2+y^2=4$
\item[(C)] $4x+y^2=4$
\item[(D)] $x+4y^2=1$
\end{list}

\bigskip\noindent\textbf{A33.}\ Find the smallest $\theta\in[0,2\pi]$ for which $r=2\cos(5\theta)$ passes through the origin.
\begin{list}{}{\setlength{\leftmargin}{3em}\setlength{\labelwidth}{2.5em}\setlength{\itemsep}{0pt}\setlength{\topsep}{2pt}}
\item[(A)] $0$
\item[(B)] $\pi/20$
\item[(C)] $\pi/10$
\item[(D)] $\pi/5$
\end{list}

\bigskip\noindent\textbf{A34.}\ For which $\theta\in[0,\pi]$ do $r=3$ and $r=2+2\cos\theta$ intersect?
\begin{list}{}{\setlength{\leftmargin}{3em}\setlength{\labelwidth}{2.5em}\setlength{\itemsep}{0pt}\setlength{\topsep}{2pt}}
\item[(A)] $\pi/6$
\item[(B)] $\pi/4$
\item[(C)] $\pi/3$
\item[(D)] $2\pi/3$
\end{list}

\bigskip\noindent\textbf{B1.}\ The graph $f(x)=2e^{\sin x}-3$ crosses the $x$-axis once on $[0,1]$. What is the $x$-coordinate?
\begin{list}{}{\setlength{\leftmargin}{3em}\setlength{\labelwidth}{2.5em}\setlength{\itemsep}{0pt}\setlength{\topsep}{2pt}}
\item[(A)] $0.209$
\item[(B)] $0.417$
\item[(C)] $0.552$
\item[(D)] $0.891$
\end{list}

\bigskip\noindent\textbf{B2.}\ Find the $x$-intercept of $f(x)=\sqrt{\sin x+1}+x^3-3e^{\cos x}+6$ on the portion where $f$ is decreasing.
\begin{list}{}{\setlength{\leftmargin}{3em}\setlength{\labelwidth}{2.5em}\setlength{\itemsep}{0pt}\setlength{\topsep}{2pt}}
\item[(A)] $-1.334$
\item[(B)] $-0.065$
\item[(C)] $-0.801$
\item[(D)] $0.472$
\end{list}

\bigskip\noindent\textbf{B3.}\ For $f(x)=e^{x^5-2x^3+2}-7$ on $[-2,2]$, find all $x$ for which $f(x)<0$. Endpoints shown numerically are rounded.
\begin{list}{}{\setlength{\leftmargin}{3em}\setlength{\labelwidth}{2.5em}\setlength{\itemsep}{0pt}\setlength{\topsep}{2pt}}
\item[(A)] $[-2,-1.421)$ and $(0.305,1.407)$
\item[(B)] $(-1.421,0.305)$ only
\item[(C)] $[-2,-1.421)$ only
\item[(D)] $(0.305,1.407)$ only
\end{list}

\bigskip\noindent\textbf{B4.}\ How many times does $f(x)=(4-2x-2x^2)\cos(3x-4)$ cross the $x$-axis on $[-3,2]$?
\begin{list}{}{\setlength{\leftmargin}{3em}\setlength{\labelwidth}{2.5em}\setlength{\itemsep}{0pt}\setlength{\topsep}{2pt}}
\item[(A)] four
\item[(B)] five
\item[(C)] six
\item[(D)] seven
\end{list}

\bigskip\noindent\textbf{B5.}\ On $[0,2\pi]$, the curve $r=\theta-2\cos\theta$ has one point with $x$-coordinate $2$. Find its $y$-coordinate.
\begin{list}{}{\setlength{\leftmargin}{3em}\setlength{\labelwidth}{2.5em}\setlength{\itemsep}{0pt}\setlength{\topsep}{2pt}}
\item[(A)] $-4.594$
\item[(B)] $-3.764$
\item[(C)] $1.979$
\item[(D)] $5.201$
\end{list}

\bigskip\noindent\textbf{B6.}\ If $f(x)=1+e^x$, the domain of $f^{-1}$ is
\begin{list}{}{\setlength{\leftmargin}{3em}\setlength{\labelwidth}{2.5em}\setlength{\itemsep}{0pt}\setlength{\topsep}{2pt}}
\item[(A)] $(-\infty,\infty)$
\item[(B)] $(0,\infty)$
\item[(C)] $(1,\infty)$
\item[(D)] $(2,\infty)$
\end{list}

\end{document}
